\section{Conclusion and Future Work}
\label{sec:conclusion}

In this paper, we have exposed and systematically analyzed a critical vulnerability in adaptive model merging at deployment: the threat of physical test-time sensor and channel noise. 
We characterized the phenomenon of \textbf{Noise-Entropy Collapse} caused by the \emph{Overfitting-Optimizer Paradox}, demonstrating that standard test-time adaptation (AdaMerging) collapses out-of-distribution accuracy by fitting high-frequency transductive noise. 
To resolve this, we introduced \textbf{FlatMerge}, a robust, plug-and-play dual-regularization framework that: (i) restricts task-blending coefficients to a low-degree polynomial subspace of normalized layer depth, and (ii) optimizes these parameters using flatness-aware minimization in the highly compressed coefficient space. 
Our extensive empirical evaluations across two simulated environments and 15 independent random seeds demonstrate that FlatMerge achieves state-of-the-art robustness, outperforming unconstrained AdaMerging and RegCalMerge, while simultaneously reducing seed variance by more than $60\%$. 
Importantly, because our optimization space is extremely compact, FlatMerge requires **zero activation memory caching** during adaptation, satisfying the rigid hardware constraints of edge deployments.

\subsection{Limitations}
Despite its strong mathematical and empirical advantages, we identify three primary limitations of FlatMerge:
\begin{enumerate}
    \item \textbf{Choice of Polynomial Degree ($d$):} While a quadratic polynomial ($d=2$) provides an optimal regularizing subspace for a 12-layer Vision Transformer, scaling to significantly deeper architectures—such as 80-layer Large Language Models (LLMs)—may require higher-degree polynomials or piece-wise polynomial splines to capture complex hierarchical task transitions. Formally, for a network divided into $S$ structural segments with layer boundaries $0 = b_0 < b_1 < \dots < b_S = L$, we can define a piece-wise spline parameterization $\lambda^l_k = \sum_{j=0}^{d} w_{k, s, j} \left(\frac{l - b_{s-1}}{b_s - b_{s-1}}\right)^j$ for $l \in (b_{s-1}, b_s]$, subject to $C^0$ or $C^1$ continuity constraints at the segment boundaries $b_s$. This preserves a highly compact coefficient space while maintaining representational capacity in ultra-deep backbones.
    \item \textbf{Non-Stationary Noise Streams:} FlatMerge assumes a stationary noise level $\gamma$ during each test-time adaptation window. In physical scenarios where environmental corruptions change dynamically and non-stationarily (e.g., passing through a tunnel or sudden weather changes), the fixed perturbation radius $\rho = 0.05$ may need to adapt dynamically to avoid under-regularization or over-smoothing.
    \item \textbf{Weight-Reconstruction DRAM Bandwidth:} As analyzed in Section~\ref{sec:bandwidth}, on highly memory-bandwidth-constrained edge accelerators, the repeated on-the-fly reconstruction of merged weights twice per step creates a DRAM-to-SRAM transfer bottleneck ($15352.34\text{ ms/step}$ vs. $7418.40\text{ ms/step}$ in weight-space TTA).
\end{enumerate}

\subsection{Future Directions}
These limitations point to several highly promising paths for future research. 
First, we plan to implement FlatMerge inside a fused custom CUDA kernel that performs the task vector scaling and addition on-the-fly during SRAM load, completely bypassing the need for intermediate DRAM weight writes. 
Second, we will explore dynamic, low-precision quantized task-vector storage (e.g., INT8 or INT4) to reduce DRAM transactions by up to $4\times$. 
Third, we intend to study \emph{partial layer-wise adaptation}, restricting the polynomial coefficient optimization to a subset of late, task-specific layers, which would reduce the number of reconstructed parameters and completely eliminate the latency bottleneck while preserving FlatMerge's exceptional out-of-distribution robustness.
Finally, we plan to calibrate and evaluate our continuous simulation sandbox on high-dimensional datasets (such as ImageNet-1K task subsets and multi-task text classification benchmarks) and release our complete simulation and profiling code as an open-source library to facilitate rapid, low-compute prototyping of test-time model merging algorithms by the broader community.
